\subsection{Shapeev (2016): Moment Tensor Potentials}
\label{sec:appendix-example-shapeev2016}

\paragraph{Q1 -- Bibliographic identification.}
Shapeev introduces the Moment Tensor Potential (MTP), a class
of systematically improvable nonparametric interatomic
potentials based on permutation-, rotation-, and
reflection-invariant polynomials of moment tensors of the local
atomic environment. The numerical experiments fit MTP to the
publicly available tungsten DFT database of
Szlachta--Bart\'ok--Cs\'anyi (\texttt{libatoms.org}) and
benchmark accuracy and CPU cost against GAP. Published in
\emph{Multiscale Model.\ Simul.} 14(3), 1153--1173 (2016);
arXiv:1512.06054.
\lstinputlisting[language=turtle]{artifacts/kg/listings/shapeev2016/q01-bibliographic.ttl}

\paragraph{Q2 -- Material system.}
The lead demonstration system is body-centred-cubic tungsten,
inheriting the configuration set (bulk, surfaces, defects, and
finite-temperature MD snapshots) from the published GAP-tungsten
database.
\lstinputlisting[language=turtle]{artifacts/kg/listings/shapeev2016/q02-material.ttl}

\paragraph{Q3 -- Reference calculation method.}
Reference data are produced from Kohn--Sham DFT calculations
(inherited from the libatoms tungsten database).
\lstinputlisting[language=turtle]{artifacts/kg/listings/shapeev2016/q03-reference-method-type.ttl}

\paragraph{Q4 -- Reference settings.}
The DFT settings inherited from the libatoms tungsten database
are CASTEP with the PBE/GGA exchange--correlation functional and
norm-conserving pseudopotentials, with k-point meshes converged
per-cell and finite-temperature electronic broadening at 1000~K.
The plane-wave energy cutoff used for the database is not
restated in this methodological paper, so we omit
$\dlRole{energyCutoff}$ here and flag this in Q12.
\lstinputlisting[language=turtle]{artifacts/kg/listings/shapeev2016/q04-reference-settings.ttl}

\paragraph{Q5 -- Training dataset.}
The fit uses 9{,}693 configurations of tungsten (with nearly
150{,}000 individual atomic environments) from the libatoms
tungsten database. Energies and forces are covered.
\lstinputlisting[language=turtle]{artifacts/kg/listings/shapeev2016/q05-dataset.ttl}

\paragraph{Q6 -- Sampling strategies.}
The libatoms tungsten database combines bulk and surface
enumeration, point-defect and dislocation sampling, and
DFT MD snapshots. We record three sampling instances reflecting
those classes; the present paper does not author the database
and hence does not contribute new sampling strategies.
\lstinputlisting[language=turtle]{artifacts/kg/listings/shapeev2016/q06-sampling.ttl}

\paragraph{Q7 -- MLIP method, components, implementation.}
The method is the original MTP: total energy as a sum of local
contributions $V(Dx_k)$, with $V$ expanded as a linear
combination of moment-tensor basis functions $B_\alpha(u)$ that
are permutation-, rotation-, and reflection-invariant
polynomials of moment tensors $M_{\mu,\nu}$. The loss is a
regularised least-squares loss on energies and forces, with
$\ell_2$ or $\ell_0$ regularisation; the regularisation parameter
$\gamma$ is selected by 16-fold cross-validation. The
implementation used for the GAP comparison is the QUIP MTP
prototype.
\lstinputlisting[language=turtle]{artifacts/kg/listings/shapeev2016/q07-method.ttl}

\paragraph{Q8 -- Hyperparameter settings.}
Reported hyperparameter settings: cutoff radius
$R_{\mathrm{cut}} = 4.9$~\AA, minimal distance
$R_{\min} = 1.9$~\AA, and two MTP variants---MTP$_1$
($\deg(B_\alpha) + 8(\#\alpha) \leq 62$, $\#\alpha \leq 4$,
$\mu \leq 5$, $\nu \leq 4$, 11{,}133 basis functions) and
MTP$_2$ ($\deg(B_\alpha) + 8(\#\alpha) \leq 52$, $\#\alpha \leq 5$,
$\mu \leq 3$, $\nu \leq 5$; 760 basis functions extracted via
the $\ell_0$ algorithm with $N_{\mathrm{cap}} = 4$).
\lstinputlisting[language=turtle]{artifacts/kg/listings/shapeev2016/q08-hyperparameter-settings.ttl}

\paragraph{Q9 -- MLIP run and trained model.}
A single training run produces the production tungsten MTP,
with two basis-set variants (MTP$_1$ and MTP$_2$) recorded as
hyperparameter settings on the same run rather than as separate
runs/models.
\lstinputlisting[language=turtle]{artifacts/kg/listings/shapeev2016/q09-run-and-model.ttl}

\paragraph{Q10 -- Benchmark results.}
Headline accuracy: MTP$_1$ attains a force RMSE of
0.0427~eV/\AA{} on the full training set (relative RMS
2.8\%) and 0.0511~eV/\AA{} under 16-fold cross-validation
(relative RMS 3.4\%) versus DFT-PBE; on the same database GAP
attains 0.0633~eV/\AA{} (4.2\%). MTP$_2$ matches GAP's accuracy
with $\sim 13\times$ fewer parameters.
\lstinputlisting[language=turtle]{artifacts/kg/listings/shapeev2016/q10-benchmark-results.ttl}

\paragraph{Q11 -- Computational resources.}
Per-atom inference time is reported in Table~1: 2.9~ms/atom
for MTP$_1$, 0.8~ms/atom for MTP$_2$, and 134.2~ms/atom for the
GAP baseline, on a single core of an Intel i7-2675QM laptop
CPU. We encode MTP$_1$'s 2.9~ms/atom as the headline
$\dlRole{inferenceTimePerAtom}$ figure with a free-text
$\dlRole{inferenceHardware}$. Training duration, total CPU
hours, and peak memory are not reported.
\lstinputlisting[language=turtle]{artifacts/kg/listings/shapeev2016/q11-resources.ttl}

\paragraph{Q12 -- Gaps.}
Beyond the partial Q11 reporting (training-cost metadata
absent; only per-atom inference time is given), the following
ontology-expressible items are not reported in this paper
itself: the DFT plane-wave $\dlRole{energyCutoff}$ used for the
inherited tungsten database; explicit
$\dlRole{frozenCore}$ treatment beyond the pseudopotential
partition; software versioning of the QUIP MTP prototype; an
explicit named individual for the cross-validation training
algorithm. The protocol assumes one method $\times$ one
material per file; the present paper is methodological, so the
tungsten benchmark is encoded as the lead system and the
absence of additional materials is itself a deliberate
simplification.
